\scanpage{138}
\mspage{87}

$t'$ est l'unique élément de $T'(\simeq T)\unsure{,}=X'\setminus U'$. Donc
il faut des conditions sur $U_{/t}$ \unsure{(?)} \deleted{\unsure{même}} qui
assurent que non seulement $U_{/t}$, mais aussi
toutes les catégories similaires $U'_{/t}$, soient
$\mathcal W$-asphériques. Dans le cas général \deleted{\unsure{de même}}\inserted{où}
$T$ est quelconque, c'est cette condition là que
nous allons imposer aux $U_{/t}$ $(t\in\mathrm{Ob}\,T)$, et il
s'avérera qu'\deleted{elles \unsure{sont}} cette fois-ci, \unsure{elles|elle} seront stables
par \inserted{tout} \unsure{cochangement|changement} de base $U\xrightarrow{p}U'$, i.e. que les
$U'_{/t'}$ $(t'\in\mathrm{Ob}\,T')$ satisferont alors à la \unsure{=|même} condition,
\inserted{C'est cette condition} que nous allons à présent reformuler.

2.3. \uline{Objets totalement $\mathcal W$-asphériques de} $\mathrm{Cat}$, \uline{et
morphismes totalement $\mathcal W$-asphériques.}

Reprenons la situation
\diagram{\begin{tikzcd}[ampersand replacement=\&]
 U \arrow[r,hook,"i"] \& X \& T \arrow[l,hook',"j"']
\end{tikzcd}}
, avec $T=\{t\}$, catégorie
finale \inserted{$t\in\overline U$, i.e. $X=\overline U$,} $e$,
et celle d'un $\mathrm{cochg}^{\mathrm t}$ de base $U\unsure{-|\to}U'$,
donnant naissance à un carré cocartésien dans $\mathrm{Cat}$ :
\diagram{\((\unsure{\alpha|1|*})\qquad\begin{tikzcd}[ampersand replacement=\&]
 U \arrow[r,hook,"i"] \arrow[d,"p"'] \& X \arrow[d,"q"'] \& T=\{t\} \arrow[l,dotted,hook',"j"'] \arrow[d,dotted,"\mathrm{iso}"] \\
 U' \arrow[r,hook,"i'"] \& X' \& T'=\{t'\} \arrow[l,dotted,hook',"j'"']
\end{tikzcd}\)}.
\deleted{Fixons nous \unsure{sur} le} \deleted{\unsure{Le|=} calcul de} La donnée de recollement $H$
entre $U$ et $T$, $H:\underbrace{U^{\circ}\times T}_{\simeq U^\circ}-\mathrm{Ens}$, \deleted{\unsure{peut}}
s'interprète comme un préfaisceau sur $U$,
\inserted{$H\in\mathrm{Ob}\,U^\wedge\ \unsure{{:}{:}|{:}}\ U^\circ\to\mathrm{Ens}$}

% archive sheet number: 138
